\unitopen{45}
\begin{candos}
\item prowadzić small talk w biurze i na \de{Apéro} (pogoda, weekend, praca, podróże);
\item przekazać krytykę i feedback uprzejmie, rzeczowo i bez oskarżeń;
\item rozwiązywać nieporozumienia i konflikty w zespole;
\item dopasować rejestr (\de{du}/\de{Sie}, formalny/nieformalny) do sytuacji.
\end{candos}

\section{Słownictwo}
\vocabtable{45}

\section{Small talk}
Dobre tematy: pogoda, weekend, podróże, jedzenie, sport, praca (ogólnie), miejsce urodzenia, miasto. Unikaj: pieniędzy, polityki, religii,
szczegółów prywatnych. Zagajenia: \de{Wie war Ihr Wochenende?} \de{Schönes Wetter heute, oder?} \de{Waren Sie schon einmal in Polen?}
Podtrzymywanie: \de{Ach, wirklich?} \de{Das klingt interessant.} \de{Und wie war es?} Zakończenie: \de{Es war schön, mit Ihnen zu sprechen.}
\de{Ich muss leider weiter, aber wir sehen uns bestimmt.}

\section{Feedback i konflikt}
\begin{center}\small
\begin{tabular}{@{}p{0.27\linewidth}p{0.65\linewidth}@{}}
\toprule
\textsf{\textbf{Cel}} & \textsf{\textbf{Wzór}}\\
\midrule
Wprowadzenie & \de{Darf ich kurz etwas ansprechen?}\\
Obserwacja & \de{Mir ist aufgefallen, dass der Report heute zum dritten Mal zu spät war.}\\
Wpływ & \de{Das führt dazu, dass wir die Daten nicht rechtzeitig prüfen können.}\\
Prośba & \de{Könnten wir vereinbaren, dass er bis zehn Uhr vorliegt?}\\
Słuchanie & \de{Wie sehen Sie das?}\\
Zbieżność & \de{Ich verstehe Ihren Standpunkt, aber \dots}\\
\bottomrule
\end{tabular}
\end{center}

\begin{zasada}[Rejestr]
\de{Sie} i formalność w kontakcie z klientem i na początku. \de{du} w zespołach, szczególnie w IT. Przełącz się na \de{du} dopiero po propozycji
rozmówcy. Zmiana rejestru bez zgody jest niegrzeczna, tak jak i trzymanie się \de{Sie} po prośbie o \de{du}.
\end{zasada}

\begin{falle}
Polska bezpośredniość („to jest źle”) w Szwajcarii bywa odbierana jako agresja. Zamiast \de{Das ist falsch} mów \de{Ich sehe das etwas anders}
albo \de{Mir ist aufgefallen, dass \dots}. Drugi błąd: milczenie przy problemie: Szwajcarzy cenią wcześniejsze, uprzejme zasygnalizowanie.
Trzeci: używanie „wy” jako \de{ihr} wobec nieznajomych (\de{Sie}).
\end{falle}

\begin{szwajcaria}
Kultura szwajcarska ceni umiar, punktualność, konsensus i dyskrecję; konflikt rozwiązuje się rozmową w cztery oczy, nie przy zespole.
\de{Apéro} to nieformalne spotkanie z drinkiem; rozmowa jest lekka, a wyjście po godzinie jest normalne. Rozmowa o zarobkach nie jest dobrym tematem.
Dialekt zmienia się z miastem; nie udawaj, że go mówisz, ale okaż zainteresowanie: \de{Wie sagt man das auf Schweizerdeutsch?}
\end{szwajcaria}

\begin{praca}
Prośba o feedback: \de{Hätten Sie kurz Zeit für Feedback zu meiner Präsentation?} Odpowiedź na krytykę: \de{Danke für die Rückmeldung. Das
nehme ich gern auf.} Potwierdzenie nieporozumienia: \de{Ich glaube, wir haben uns missverstanden. Lassen Sie mich das noch einmal erklären.}
\end{praca}

\section{Dialogi}
\textbf{Dialog 1 --- nieporozumienie w zespole}
\begin{dialog}
\turn{A}{Darf ich kurz etwas ansprechen? Mir ist aufgefallen, dass der Report heute später kam als vereinbart.}{Czy mogę coś krótko poruszyć? Zauważyłem, że raport dziś przyszedł później niż uzgodniono.}
\turn{B}{Ach, das tut mir leid. Ich dachte, die Frist sei Freitag.}{Ach, przykro mi. Myślałem, że termin to piątek.}
\turn{A}{Dann haben wir uns wohl missverstanden. Könnten wir künftig schriftlich festhalten, bis wann was fertig sein soll?}{Wobec tego chyba źle się zrozumieliśmy. Czy moglibyśmy w przyszłości zapisywać na piśmie, do kiedy co ma być gotowe?}
\turn{B}{Gern. Ich schicke Ihnen nach dem Meeting eine kurze Zusammenfassung.}{Chętnie. Po spotkaniu wyślę Panu krótkie podsumowanie.}
\end{dialog}

\section{Ćwiczenia}
\begin{cwiczenia}
\exitem{Zacznij small talk z kolegą w poniedziałek.}{\de{Wie war dein Wochenende?}}
\exitem{Podtrzymaj rozmowę (jedna reakcja).}{\de{Ach, wirklich? Das klingt interessant.}}
\exitem{Zakończ rozmowę uprzejmie.}{\de{Ich muss leider weiter, aber wir sehen uns bestimmt.}}
\exitem{Zacznij trudną rozmowę.}{\de{Darf ich kurz etwas ansprechen?}}
\exitem{Opisz obserwację bez oceny osoby.}{\de{Mir ist aufgefallen, dass der Report zu spät kam.}}
\exitem{Poproś o uzgodnienie na przyszłość.}{\de{Könnten wir vereinbaren, dass er bis zehn Uhr vorliegt?}}
\exitem{Wyraź zrozumienie i zastrzeżenie.}{\de{Ich verstehe Ihren Standpunkt, aber ich sehe das etwas anders.}}
\exitem{Przeformułuj łagodniej: \de{Das ist falsch.}}{\de{Ich sehe das etwas anders.} lub \de{Ich bin nicht sicher, ob das stimmt.}}
\exitem{Który rejestr: do klienta, którego nie znasz?}{\de{Sie}}
\exitem{Przetłumacz: „Chyba źle się zrozumieliśmy.”}{\de{Ich glaube, wir haben uns missverstanden.}}
\end{cwiczenia}

\weekend{Odegraj z AI cztery scenki: small talk na \de{Apéro}, informacja zwrotna dla kolegi, nieporozumienie z klientem,
łagodna odmowa. Nagraj wszystkie i oceń rejestr (\de{du}/\de{Sie}) oraz uprzejmość.}
